%% kapa: exam=1 round=27 year=2016 period=1 form=A subject=economics subject_ko=경제학원론 seq=18 number=58
%% source: Q-Net 감정평가사 기출문제 (cst00309) artlSeq=5169402 file=extracted/1008/5169402/제27회 감정평가사 1차 문제/제27회 감정평가사 1교시문제(A형).pdf
%% conversion: kordoc 4.24.51 -> tools/gen1.py (auto); review=proofread-sonnet; flags=none
\begin{question}{58}
甲은 X재와 Y재 두 재화를 1 : 1 비율로 묶어서 소비한다. X재의 가격과 수요량을 각각 $P_X$와 $Q_X$라 한다. 소득이 1,000이고 Y재의 가격이 10일 때 甲의 X재 수요함수로 옳은 것은? (단, 소비자는 효용을 극대화하고 소득을 X재와 Y재 소비에 모두 지출한다.)
\begin{choices}
  \item $Q_X = 1{,}000/(10 + P_X)$
  \item $Q_X = 990 - P_X$
  \item $Q_X = 500 - P_X$
  \item $Q_X = 1{,}000 - P_X$
  \item $Q_X = 500/P_X$
\end{choices}
\end{question}
\answer{1}
